\section{Analytical Framework}
\label{sec:framework}

The problem setting above specifies what CUAs observe and how they act, but a survey still needs a synthesis device: a way to compare systems, risks, and controls without collapsing them into platform-specific case lists. To answer that need, this paper adopts a two-dimensional framework that combines an \textbf{architectural view} with a \textbf{lifecycle view}. The architectural view explains how a CUA is assembled; the lifecycle view explains when its capabilities, risks, and controls are formed, exposed, stressed, and revised.

\subsection{Why a Joint View Is Needed}
For survey purposes, architecture alone is necessary but incomplete. A component-level description can explain how a CUA perceives the environment, reasons over a task, invokes tools, and executes interface operations. That view is useful for comparing agent designs across web, desktop, and mobile settings.

However, many important failures cannot be explained from structure alone. A weak perceptual module may reflect poor training data, but the same downstream error may also be caused by deployment-time observation loss, runtime interface deception, or maintenance-stage drift. Likewise, the same execution interface may be relatively benign under strict sandboxing but high-risk under broad permissions. A purely structural account therefore tells us \emph{how} the system works, but not \emph{when} and \emph{under what conditions} it becomes unsafe or hard to control.

Lifecycle analysis provides that missing dimension. We distinguish four stages:
\emph{Creation}, \emph{Deployment}, \emph{Operation}, and \emph{Maintenance}. Creation shapes what the agent can learn from data, objectives, and action abstractions. Deployment binds those capabilities to real platforms, tools, permissions, and trust boundaries. Operation reveals whether the system can remain coherent under long-horizon interaction and adversarial content. Maintenance determines whether the system stays reliable as interfaces, workflows, tools, and threat surfaces evolve. This view helps distinguish where a failure originates from where it becomes visible.

\begin{table*}[!t]
\centering
\small
\caption{How the lifecycle stages concentrate pressure on the tri-layer architecture. The joint view links capability formation, failure emergence, and control placement.}
\label{tab:joint_framework}
\renewcommand{\arraystretch}{1.15}
\setlength{\tabcolsep}{4.5pt}
\begin{tabularx}{\textwidth}{L{1.7cm} L{2.8cm} L{2.8cm} L{2.8cm} Y}
\toprule
\textbf{Stage} & \textbf{Perception} & \textbf{Decision} & \textbf{Execution} & \textbf{Primary control question} \\
\midrule
Creation & Learn actionable grounding from data and representation choices & Shape objectives, decomposition priors, and reasoning habits & Define action vocabulary and safety abstractions & What capabilities and unsafe priors are built into the agent before deployment? \\
Deployment & Bind observation channels such as pixels, DOM, OCR, or middleware state & Attach orchestration, memory, and policy mediation to the task loop & Bind tools, privileges, and sandboxes to executable actions & What authority is exposed, and how is model intent translated into controlled execution? \\
Operation & Handle ambiguity, visual deception, and missing state online & Sustain long-horizon coherence, memory, recovery, and constraint tracking & Execute under latency, TOCTOU, and real permission boundaries & How do local errors or adversarial inputs become unsafe trajectories in live use? \\
Maintenance & Adapt to UI drift, localization regressions, and benchmark shift & Update policies, evaluators, and recovery logic without regressions & Patch tools, permissions, logs, and containment mechanisms & How is reliability preserved as the environment and threat model change? \\
\bottomrule
\end{tabularx}
\end{table*}

\subsection{What the Joint Framework Explains}
The joint framework does more than organize prior work. It supports three analytical claims that guide the remainder of the survey.

\textbf{First, capability is lifecycle-dependent.} What appears at runtime as ``agent ability'' is produced jointly by representational choices, training signals, deployment-time bindings, and online recovery mechanisms. A strong model alone is therefore not a sufficient explanation for robust computer use.

\textbf{Second, failure origin and failure manifestation are different questions.} Many failures surface during operation, but their roots may lie in creation-stage data gaps, deployment-stage trust decisions, or maintenance-stage regressions. Treating every unsafe action as a runtime reasoning failure obscures the real intervention point.

\textbf{Third, control placement must follow both architecture and lifecycle.} Perception problems call for different controls than execution problems, and creation-stage controls differ from deployment-stage or maintenance-stage controls. Security and governance are therefore not post hoc add-ons; they are constraints that must be attached to the right layer at the right time.

Table~\ref{tab:joint_framework} is the paper's master synthesis device. Later tables should therefore be read as \emph{local instantiations} of this joint view rather than as separate taxonomies. The survey's main argument is singular: CUAs should be compared through the interaction between structural layer and lifecycle stage.

\subsection{OpenClaw as a Running Example of Open Deployment}
To keep the framework concrete, we use OpenClaw as a running example throughout the remainder of the paper. Public descriptions suggest a local-first personal assistant that can be reached through persistent communication channels and connected to longer-lived sessions, tools, and agent communities~\cite{openclaw2026website,ibm2026openclaw,chen2026openclaw}. These publicly described properties make it analytically useful because they expose, in a single setting, the main system tensions emphasized by this paper: how perception, decision, and execution are mediated through an open control plane; how deployment choices concentrate authority; and how maintenance and governance become harder once tools and capabilities can evolve or propagate socially.

This use of OpenClaw is illustrative rather than exhaustive. We do not assume access to its full internal design, nor do we treat it as the only valid deployment model for CUAs. Instead, OpenClaw serves as a concrete example of the broader class of \emph{open-deployment CUA systems}, in which persistent ingress channels, extensible tool interfaces, memory, and ecosystem participation produce capabilities and risks that benchmark-only analyses tend to understate.

\subsection{Orientation for the Rest of the Paper}
The remainder of the survey proceeds in two passes. We first analyze the tri-layer architecture of CUAs and then examine the four lifecycle stages through which those systems are built, deployed, operated, and maintained. Later sections use the joint framework to organize the security landscape and derive practical implications for defense and governance~\cite{kuntz2025harm,zheng2026riskybench,wang2026asset}.
